% Plantilla mínima de los boletines IES (pandoc --template).
% Solo carga lo estrictamente necesario por el contenido; la tipografía y la
% geometría las decide LaTeX (tamaño carta, 12pt, márgenes por defecto).

\documentclass[12pt,letterpaper]{article}

% Codificación y acentos (pdflatex)
\usepackage[T1]{fontenc}
\usepackage[utf8]{inputenc}

% Fuente escalable (Latin Modern); evita las bitmap EC de METAFONT
\usepackage{lmodern}

% Español: títulos y pies ("Figura"), guionado
\usepackage[spanish]{babel}

% Secciones sin numeración automática (equivalente a \section*{...}): solo se
% muestran los números que el markdown escribe a mano ("1.", "2.", "A.", ...).
\setcounter{secnumdepth}{-1}

% Los títulos nunca quedan huérfanos al pie de página: reserva espacio antes
% de cada sección y, si no cabe el título + unas líneas, salta a la página
% siguiente. Necesario sobre todo cuando el título precede a una tabla longtable.
\usepackage{needspace}
\usepackage{etoolbox}
\pretocmd{\section}{\needspace{3\baselineskip}}{}{}
\pretocmd{\subsection}{\needspace{3\baselineskip}}{}{}
\pretocmd{\subsubsection}{\needspace{3\baselineskip}}{}{}

% Párrafos: penalizar líneas huérfanas (club) y viudas (widow) al cambiar de
% página: la primera línea de un párrafo no queda sola abajo, ni la última sola arriba.
\clubpenalty=10000
\widowpenalty=10000
\displaywidowpenalty=10000

% Matemáticas (variables como T, I, neq)
\usepackage{amsmath,amssymb}

% Imágenes: los diagramas PDF se ajustan al ancho de caja, sin anchos fijos.
% float + [H] coloca cada figura donde el Markdown la sitúa: evita páginas
% flotantes con una sola figura y títulos separados de su diagrama.
\usepackage{graphicx}
\usepackage{float}
\floatplacement{figure}{H}
\makeatletter
\def\maxwidth{\ifdim\Gin@nat@width>\linewidth\linewidth\else\Gin@nat@width\fi}
\def\maxheight{\ifdim\Gin@nat@height>\textheight\textheight\else\Gin@nat@height\fi}
\makeatother
\setkeys{Gin}{width=\maxwidth,height=\maxheight,keepaspectratio}

% Fondo: membrete a página completa (carta) en TODAS las hojas. La ruta la
% inyecta build-boletin.sh con -V membrete=<ruta absoluta>; si la variable no
% se define, el fondo simplemente no se carga. Ojo: en un template de pandoc
% los signos de dolar se interpretan como variables de template, incluso en los
% comentarios; por eso aqui no debe aparecer ningun signo de dolar suelto
% aparte de la variable membrete del bloque condicional.
\usepackage{background}
% NOTA: hay que dar width Y height explicitos (paperwidth x paperheight). Las
% claves globales de Gin (width=\maxwidth, height=\maxheight, keepaspectratio)
% fijadas mas arriba para los diagramas se aplican por defecto a todo
% \includegraphics; con solo width=\paperwidth el fondo quedaria reducido al
% area de texto (margenes). Con ambas dimensiones explicitas y la misma
% proporcion carta del membrete, llena la pagina sin deformar ni recortar.
\backgroundsetup{
  scale=1,
  angle=0,
  opacity=1,
  contents={\includegraphics[width=\paperwidth,height=\paperheight]{/home/kritt/Git/boletin/.opencode/skills/boletin-latex-pdf/assets/membrete.pdf}}
}

% Tablas (los boletines pueden traer tablas Markdown)
\usepackage{longtable,booktabs,array,calc}

% Pies de figura
\usepackage{caption}
\captionsetup{font=small,labelfont=bf}

% Listas compactas de Markdown
\providecommand{\tightlist}{%
  \setlength{\itemsep}{0pt}\setlength{\parskip}{0pt}}

% Enlaces y destinos (al final, como recomienda hyperref)
\usepackage{hyperref}

\begin{document}

\hypertarget{termodinuxe1mica-financiera-09-duxe9ficit-comercial-muxe9xico-ee.-uu.}{%
\section{Termodinámica Financiera 09: Déficit Comercial México-EE.
UU.}\label{termodinuxe1mica-financiera-09-duxe9ficit-comercial-muxe9xico-ee.-uu.}}

\href{https://www.eleconomista.com.mx/empresas/aumenta-deficit-comercial-estados-unidos-mexico-septiembre-20261109-842947.html}{Déficit
comercial México-EE. UU.} ·
\href{https://www.iesencial.com/boletin/\#art=2026-11-12-Boletin-09\%3A\%20Deficit\%20comercial\%20mex-usa.md}{Publicado
en iesencial.com}

\hypertarget{titular-y-resumen-ejecutivo}{%
\subsection{1. Titular y Resumen
Ejecutivo}\label{titular-y-resumen-ejecutivo}}

\textbf{Aumenta 26.1\% el déficit comercial de Estados Unidos con México
a septiembre de 2026: La trampa sistémica del dólar, la fragilidad del
nearshoring y la ilusión de la eficiencia local.}

Mientras la política económica estadounidense interpreta el déficit
bilateral de \$184,019M (con exportaciones mexicanas de \$481,586M) como
una hemorragia industrial que exige aranceles, el análisis sistémico
revela una realidad opuesta: Estados Unidos opera atrapado en el
\textbf{Dilema de Triffin}, exportando dólares para sostener su
hegemonía financiera a costa de su capacidad física. Al integrar la
\textbf{Teoría de Restricciones (TOC)} con el \textbf{Efecto Séneca} de
Ugo Bardi, demostramos que restringir este flujo mediante barreras
comerciales no reindustrializa a Norteamérica, sino que empuja a la
cadena de suministro regional hacia un precipicio de colapso no lineal.

\begin{center}\rule{0.5\linewidth}{0.5pt}\end{center}

\hypertarget{marco-teuxf3rico-la-truxedada-sistuxe9mica-del-colapso-y-el-enfoque}{%
\subsection{2. Marco Teórico: La Tríada Sistémica del Colapso y el
Enfoque}\label{marco-teuxf3rico-la-truxedada-sistuxe9mica-del-colapso-y-el-enfoque}}

Para comprender la verdadera naturaleza del desequilibrio comercial
entre México y Estados Unidos, integramos tres lentes fundamentales bajo
la arquitectura lógica de la \textbf{Teoría de Restricciones (TOC)}:

\begin{enumerate}
\def\labelenumi{\arabic{enumi}.}
\item
  \textbf{Robert Triffin (La Contradicción Monetaria Estructural):} El
  Dilema de Triffin establece que el país emisor de la moneda de reserva
  global (el dólar estadounidense) está obligado a proveer liquidez al
  sistema mundial incurriendo en déficits comerciales crónicos. Sin
  embargo, la acumulación perpetua de pasivos externos erosiona su base
  manufacturera interna y la confianza en su solvencia a largo plazo.
  Bajo la TOC, el diseño monetario global fuerza a que la liquidez
  financiera opere estructuralmente en contra de la capacidad física de
  producción local (\(I\) y \(OE\)).
\item
  \textbf{Ugo Bardi (La Dinámica Temporal y el Efecto Séneca):} El
  Efecto Séneca demuestra que los sistemas complejos sometidos a estrés
  acumulado y rendimientos decrecientes no declinan de forma lineal o
  gradual. Mientras el crecimiento y la integración comercial son
  procesos lentos y laboriosos, el colapso ante un choque exógeno
  ---como una guerra arancelaria o una ruptura logística--- es abrupto,
  asimétrico y acelerado (``el crecimiento es lento, pero la ruina es
  rápida'').
\item
  \textbf{Eliyahu M. Goldratt (El Eje de Control de Sistemas
  Complejos):} La Teoría de Restricciones trasciende la mera
  optimización de valor para erigirse como la ley fundamental de la
  dinámica de sistemas complejos: \textbf{el desempeño global de
  cualquier sistema interdependiente no está determinado por la suma de
  sus eficiencias locales, sino que está gobernado y controlado por un
  único elemento (la restricción o cuello de botella)}. Intentar
  optimizar subcomponentes aislados ---como forzar el equilibrio de la
  balanza bilateral--- sin elevar la capacidad de la restricción real
  del sistema destruye el rendimiento total y precipita la caída en la
  curva de Séneca.
\end{enumerate}

\begin{center}\rule{0.5\linewidth}{0.5pt}\end{center}

\hypertarget{diagnuxf3stico-del-sistema-seguxfan-toc}{%
\subsection{3. Diagnóstico del Sistema según
TOC}\label{diagnuxf3stico-del-sistema-seguxfan-toc}}

El sistema bajo análisis es el \textbf{bloque industrial y monetario de
Norteamérica}. El flujo real revela que México no actúa como un
extractor de riqueza de Estados Unidos, sino como la válvula de escape y
extensión productiva que permite a las corporaciones norteamericanas
sostener su oferta ante la escasez de capacidad manufacturera interna.

\hypertarget{identificaciuxf3n-de-la-restricciuxf3n-y-paruxe1metros-t-i-oe}{%
\subsubsection{\texorpdfstring{Identificación de la Restricción y
Parámetros (\(T\), \(I\),
\(OE\))}{Identificación de la Restricción y Parámetros (T, I, OE)}}\label{identificaciuxf3n-de-la-restricciuxf3n-y-paruxe1metros-t-i-oe}}

\begin{itemize}
\tightlist
\item
  \textbf{La Restricción del Sistema (Cuello de Botella):} La
  \textbf{capacidad física, energética y logística integrada de
  Norteamérica} (infraestructura transfronteriza, aduanas, generación
  eléctrica y mano de obra técnica calificada) para transformar insumos
  en bienes finales competitivos frente a los bloques asiáticos.
\item
  \textbf{Throughput (\(T\)):} La tasa a la que el bloque norteamericano
  genera dinero fresco mediante ventas reales al consumidor final. El
  incremento de 20.7\% en las importaciones estadounidenses de bienes
  mexicanos ---con el equipo de cómputo como principal impulsor---
  representa el flujo vital que mantiene vivo el Throughput de las
  empresas estadounidenses.
\item
  \textbf{Inversión / Inventario (\(I\)):} Todo el capital inmovilizado
  en plantas industriales, relocalización de cadenas de suministro
  (nearshoring), inventarios en tránsito y los propios bonos del Tesoro
  estadounidense acumulados en el exterior como contrapartida del
  déficit comercial (Dilema de Triffin).
\item
  \textbf{Gasto de Operación (\(OE\)):} El costo sistémico necesario
  para convertir el inventario en Throughput. Los aranceles, las
  fricciones aduanales y la duplicación de cadenas productivas inflan
  artificialmente el \(OE\) sin añadir un solo dólar de valor al
  sistema.
\end{itemize}

\hypertarget{procesos-de-pensamiento-el-dilema-sistuxe9mico-de-norteamuxe9rica}{%
\subsubsection{Procesos de Pensamiento: El Dilema Sistémico de
Norteamérica}\label{procesos-de-pensamiento-el-dilema-sistuxe9mico-de-norteamuxe9rica}}

\begin{figure}
\centering
\includegraphics{Diagramas/figura-1.pdf}
\caption{Nube de Evaporación: El Dilema de Triffin y el Déficit
Bilateral en Norteamérica}
\end{figure}

\begin{center}\rule{0.5\linewidth}{0.5pt}\end{center}

\hypertarget{dicotomuxeda-clave-1-contabilidad-de-costos-tradicional-vs.-contabilidad-del-throughput-ta}{%
\subsection{4. Dicotomía Clave 1: Contabilidad de Costos Tradicional
vs.~Contabilidad del Throughput
(TA)}\label{dicotomuxeda-clave-1-contabilidad-de-costos-tradicional-vs.-contabilidad-del-throughput-ta}}

La lectura política del déficit de \$184,019M está secuestrada por los
principios obsoletos de la \textbf{Contabilidad de Costos}, que asume
que cada eslabón de una cadena debe mostrar un saldo positivo aislado
para que el conjunto sea rentable.

\begin{itemize}
\tightlist
\item
  \textbf{La Visión de la Contabilidad de Costos Tradicional:} Observa
  la balanza comercial bilateral como un estado de resultados estático.
  Concluye que importar \$481,586M desde México equivale a una ``pérdida
  neta'' de empleos y capital. Su receta es la optimización local:
  cerrar la frontera mediante aranceles para forzar que la manufactura
  regrese a territorio estadounidense, asumiendo que reducir el déficit
  bilateral mejorará automáticamente el desempeño nacional.
\item
  \textbf{La Visión de la Contabilidad del Throughput (TA):} Demuestra
  que en sistemas altamente interdependientes, \textbf{la suma de los
  óptimos locales no es igual al óptimo global}. Gran parte de las
  importaciones desde México son bienes intermedios y manufactura
  integrada cuya producción dentro de EE. UU. es físicamente inviable a
  corto plazo debido a la escasez de mano de obra y costos energéticos.
  Estrangular el flujo proveniente de México mediante aranceles eleva el
  \(OE\) y paraliza las líneas de ensamblaje finales en EE. UU.,
  destruyendo el Throughput (\(T\)) de todo el sistema y detonando el
  \textbf{Efecto Séneca}: una caída abrupta en la disponibilidad de
  bienes y un repunte inflacionario severo.
\end{itemize}

\begin{center}\rule{0.5\linewidth}{0.5pt}\end{center}

\hypertarget{dicotomuxeda-clave-2-teoruxeda-monetaria-moderna-mmt-vs.-teoruxeda-de-restricciones-toc}{%
\subsection{5. Dicotomía Clave 2: Teoría Monetaria Moderna (MMT)
vs.~Teoría de Restricciones
(TOC)}\label{dicotomuxeda-clave-2-teoruxeda-monetaria-moderna-mmt-vs.-teoruxeda-de-restricciones-toc}}

El debate sobre el déficit comercial estadounidense suele polarizarse
entre quienes temen la insolvencia fiscal y quienes proclaman la
omnipotencia del emisor monetario.

\begin{itemize}
\tightlist
\item
  \textbf{Visión de la Teoría Monetaria Moderna (MMT):} Argumenta que
  Estados Unidos, al ser el emisor monopolista del dólar en un régimen
  de tipo de cambio flotante, no enfrenta ninguna restricción financiera
  nominal. Bajo la identidad de balances sectoriales, el déficit
  comercial de EE. UU. es simplemente el superávit de ahorro en dólares
  del resto del mundo (incluido México). Para la MMT, EE. UU. obtiene
  bienes reales a cambio de asientos contables digitales, y el único
  límite real es la inflación interna.
\item
  \textbf{Visión de la Teoría de Restricciones (TOC) y el Límite
  Físico:} La TOC coincide con la MMT en que el dinero es una creación
  contable y no la restricción primaria, pero advierte su punto ciego:
  \textbf{el dinero es solo un habilitador; la restricción física y
  termodinámica siempre manda}. Inyectar liquidez o mantener el consumo
  mediante emisión monetaria sin elevar la capacidad física de la
  restricción real no genera riqueza, sino que infla el Inventario
  especulativo (\(I\)) y el ruido sistémico.
\item
  \textbf{Contraste Crítico con Triffin y Bardi:} Mientras la MMT asume
  que el mundo aceptará dólares indefinidamente sin fricciones, el
  \textbf{Dilema de Triffin} señala que la desindustrialización física
  corroe eventualmente el respaldo real de esa moneda. Cuando la
  capacidad productiva interna se atrofia, el sistema pierde
  resiliencia; si un choque arancelario o geopolítico corta el
  suministro externo (México), la emisión de dólares no puede
  materializar fábricas ni energía de la noche a la mañana, precipitando
  al sistema por el \textbf{acantilado de Séneca}.
\end{itemize}

\begin{center}\rule{0.5\linewidth}{0.5pt}\end{center}

\hypertarget{conclusiuxf3n-y-prescripciuxf3n-estratuxe9gica-los-5-pasos-de-enfoque-de-goldratt}{%
\subsection{6. Conclusión y Prescripción Estratégica: Los 5 Pasos de
Enfoque de
Goldratt}\label{conclusiuxf3n-y-prescripciuxf3n-estratuxe9gica-los-5-pasos-de-enfoque-de-goldratt}}

\emph{``Dime cómo me mides y te diré cómo me comporto.''} Si Estados
Unidos continúa midiendo su éxito económico mediante la balanza
comercial bilateral, ejecutará políticas autodestructivas que acelerarán
el colapso de su propia cadena de suministro. Para gobernar este sistema
complejo sin caer en el precipicio de Séneca, la estrategia binacional
debe ejecutarse bajo los \textbf{5 Pasos de Enfoque de Goldratt}:

\begin{enumerate}
\def\labelenumi{\arabic{enumi}.}
\tightlist
\item
  \textbf{IDENTIFICAR la restricción del sistema:} Abandonar la obsesión
  por el déficit contable bilateral. Reconocer que el verdadero cuello
  de botella de Norteamérica es su \textbf{capacidad física integrada}
  (infraestructura energética, cruces fronterizos, logística ferroviaria
  y talento técnico) para producir bienes complejos a escala global.
\item
  \textbf{EXPLOTAR la restricción:} Maximizar el rendimiento de la
  infraestructura transfronteriza existente sin esperar nuevas
  inversiones masivas. Esto implica eliminar fricciones burocráticas,
  digitalizar al 100\% los despachos aduaneros, unificar estándares
  regulatorios y garantizar suministro eléctrico ininterrumpido a los
  corredores industriales actuales del norte de México y el sur de EE.
  UU.
\item
  \textbf{SUBORDINAR todo lo demás a la decisión anterior:} Subordinar
  la política arancelaria, fiscal y monetaria al flujo continuo de la
  cadena de suministro norteamericana. Ninguna medida proteccionista
  local debe implementarse si interrumpe o encarece los insumos críticos
  que alimentan la capacidad productiva del bloque.
\item
  \textbf{ELEVAR la restricción del sistema:} Invertir de forma
  coordinada el capital atrapado en activos financieros improductivos
  (\(I\)) hacia la expansión física de la capacidad energética (redes de
  transmisión, energía limpia y gas natural), puertos de entrada, vías
  férreas y formación técnica binacional. Elevar este cuello de botella
  es la única vía real para reducir la dependencia de Asia sin destruir
  el Throughput.
\item
  \textbf{PREVENIR LA INERCIA (Volver al Paso 1):} Una vez que la
  integración logística y energética de Norteamérica deje de ser el
  cuello de botella, la restricción se desplazará hacia la innovación
  tecnológica o la disponibilidad de materias primas críticas. Las
  instituciones deben evitar que las políticas diseñadas para resolver
  la crisis comercial de 2026 se conviertan en dogmas rígidos que
  asfixien la siguiente etapa de evolución sistémica.
\end{enumerate}

\textbf{Fis. Héctor Hernández Caso}

contacto@iesencial.com

\end{document}
